\SourceLeaf{622}
\begin{center}[12.]\end{center}
\noindent$g$\quad Intensität eines galvanischen Stroms\par
\noindent$\mu$\quad Dichtigkeit des Magnetismus auf einer durch den Rheophor begrenzten Fläche\par
\noindent$\rho$\quad Abstand dieser Fläche von einer zweiten negativ magnetisirten
\[
\text{1.}\qquad g=\mu\rho
\]

2.\quad Wirkung eines electrischen Elements auf ein anderes, relativ gegen welches der Platz des erstern durch die complexe Grösse $u$ für die Zeit $t$ bestimmt wird, Entfernung $=r$
\[
-\frac{u}{r^3}-\frac{\alpha}{ur}\left(\frac{\mathrm du}{\mathrm dt}\right)^2-\frac{\beta}{r}\frac{\mathrm d\mathrm du}{\mathrm dt^2}
\]

3.\quad Gegenseitige Wirkung zwischen einem electrischen und magnetischen Elemente unter relativer Geschwindigkeit $v$
\[
\frac{\gamma v}{rr}
\]

4.\quad In einer galvanischen Strömung von der Intensität $g$, schiebt sich in der Zeit $t$ durch jeden Querschnitt die positive Electricität $\varepsilon gt$ nach der einen, die negative $-\varepsilon gt$ nach der andern Richtung.

5.\quad Es handelt sich darum die Relationen zwischen $\alpha,\beta,\gamma,\varepsilon$ zu bestimmen, $\gamma=2\beta\varepsilon$ aus der Induction bei Entstehung eines Stromes während der Zeit $\theta$ wobei die Kraft, so lange die Entstehung dauert —
\[
\begin{gathered}
=\frac{\gamma\pi hhg}{\theta RR}=\frac{2\beta\pi hh\varepsilon g}{\theta RR}\quad\text{wird}\\[6pt]
\gamma=4\alpha\varepsilon,\qquad\frac{\gamma\pi hhgv}{R^3}=\frac{4\alpha\pi hh\varepsilon gv}{R^3}\\[6pt]
2\gamma\varepsilon=1,\qquad\alpha=\frac{1}{8\varepsilon\varepsilon}=\tfrac{1}{2}\gamma\gamma,\qquad\beta=\frac{1}{4\varepsilon\varepsilon}=\gamma\gamma
\end{gathered}
\]

\begin{center}[13.]\par\medskip\textit{Grundgesetz für alle Wechselwirkungen galvanischer Ströme.}\par\smallskip(Gefunden im Juli 1835.)\end{center}
Zwei Elemente von Electricität in gegenseitiger Bewegung ziehen einander an, oder stossen einander ab, nicht eben so als wenn sie in gegenseitiger Ruhe sind.
\SourceLeaf{623}\[
\begin{gathered}
\begin{aligned}
e,x,y,z&\qquad\text{Element und Coordinaten}\\
e',x',y',z'&
\end{aligned}\\[4pt]
(x'-x)^2+(y'-y)^2+(z'-z)^2=rr
\end{gathered}
\]
Gegenseitige Wirkung (Abstossung)
\[
=\frac{ee'}{rr}\left\{1+k\left(\begin{gathered}
\left(\frac{\mathrm d(x'-x)}{\mathrm dt}\right)^2+\left(\frac{\mathrm d(y'-y)}{\mathrm dt}\right)^2\\
+\left(\frac{\mathrm d(z'-z)}{\mathrm dt}\right)^2-\tfrac{3}{2}\left(\frac{\mathrm dr}{\mathrm dt}\right)^2
\end{gathered}\right)\right\}
\]
wo $\BareRoot{\frac{1}{k}}$ eine bestimmte Geschwindigkeit vorstellt.

\begin{center}[14.]\end{center}
Auf \textit{andere} Weise stellt sich das \textit{Grundgesetz} folgendermaassen dar.

Es seien $P$ und $P'$ Punkte in zwei Strömen; $x,y,z$ und $x',y',z'$ die Coordinaten dieser Punkte, $r$ ihre Distanz; $\mathrm ds,\mathrm ds'$ zwei bei jenen Punkten anfangende Stromelemente.
\[
\begin{array}{rllll}
u&\text{Winkel}&\text{zwischen}&\mathrm ds\text{ und }\mathrm ds'\\
q&\text{-}&\text{-}&\mathrm ds\text{ und }PP'\\
q'&\text{-}&\text{-}&\mathrm ds'\text{ und }P'P
\end{array}
\]

Mit Weglassung der von der Intensität der Ströme abhangenden Factoren, üben die Elemente $\mathrm ds,\mathrm ds'$ eine gegenseitige Anziehung auf einander aus, die durch
\[
\frac{\mathrm ds\,.\,\mathrm ds'\,.\,(\cos u+\tfrac{3}{2}\cos q\,.\,\cos q')}{rr}
\]
gemessen werden kann.

Setzen wir $\mathrm dx=\mathrm dy=0$, so ist diese Kraft
\[
=\frac{\mathrm ds\,.\,\mathrm dz'}{rr}-\tfrac{3}{2}\,.\,\frac{z'-z}{r^3}\mathrm ds\,.\,\mathrm ds'\,.\,\left(\frac{\mathrm dr}{\mathrm ds'}\right)
\]
oder die partiellen Kräfte, welche $\mathrm ds$ parallel mit den Coordinatenaxen sollicitiren
\[
\begin{gathered}
\mathrm ds\,.\,\left\{\frac{(x'-x)\mathrm dz'}{r^3}-\frac{3(x'-x)(z'-z)\mathrm d'r}{2r^4}\right\}\\[4pt]
\mathrm ds\,.\,\left\{\frac{(y'-y)\mathrm dz'}{r^3}-\frac{3(y'-y)(z'-z)\mathrm d'r}{2r^4}\right\}\\[4pt]
\mathrm ds\,.\,\left\{\frac{(z'-z)\mathrm dz'}{r^3}-\frac{3(z'-z)(z'-z)\mathrm d'r}{2r^4}\right\}
\end{gathered}
\]
\SourceLeaf{624}oder, wenn man die Kräfte
\[
\begin{gathered}
-\mathrm ds\,.\,\mathrm d'\frac{(x'-x)(z'-z)}{2r^3}\\[4pt]
-\mathrm ds\,.\,\mathrm d'\frac{(y'-y)(z'-z)}{2r^3}\\[4pt]
-\mathrm ds\,.\,\mathrm d'\frac{(z'-z)(z'-z)}{2r^3}
\end{gathered}
\]
hinzusetzt, was, in sofern $\mathrm ds'$ Element eines geschlossenen Stroms ist, erlaubt ist
\[
\begin{gathered}
\frac{\mathrm ds}{2r^3}((x'-x)\mathrm dz'-(z'-z)\mathrm dx')\\[4pt]
\frac{\mathrm ds}{2r^3}((y'-y)\mathrm dz'-(z'-z)\mathrm dy')\\[4pt]
\frac{\mathrm ds}{2r^3}((z'-z)\mathrm dz'-(z'-z)\mathrm dz')
\end{gathered}
\]
Die letzte offenbar $=0$. Diesen Kräften aequivaliren aber offenbar folgende, insofern auf $\mathrm ds$ gewirkt wird
\[
\begin{aligned}
\text{1. In der Richtung }PP'\; .\;.\;.\;.\;.\quad&\frac{\mathrm ds\,.\,\mathrm dz'}{2rr}=\frac{\mathrm ds\,.\,\mathrm ds'}{2rr}\cos u\\[4pt]
\text{2. in der Richtung parallel mit }\mathrm ds'\; .\;.\;.\;.\;.\quad&-\frac{\mathrm ds\,.\,\mathrm ds'}{2rr}\,.\,\cos q
\end{aligned}
\]
welchen man noch beifügen darf
\[
\text{3. in der Richtung }\mathrm ds\; .\;.\;.\;.\;.\;.\quad+\frac{\mathrm ds\,.\,\mathrm ds'}{2rr}\,.\,\cos q'
\]
wogegen dann auf $\mathrm ds'$ drei diesen genau entgegengesetzte Kräfte wirken werden.

Sind die Coordinaten des wirkenden Stromelements $x,y,z$, die des Elements, auf welches gewirkt wird $0,0,0$; die Richtung und Stärke des ersten und zweiten Elements nach den Coordinaten geschätzt $\xi',\eta',\zeta';\ \xi,\eta,\zeta$, so ist nach \textsc{Ampère} die ganze Kraft anziehend
\[
\frac{\xi\xi'+\eta\eta'+\zeta\zeta'}{rr}-\tfrac{3}{2}\frac{(x\xi+y\eta+z\zeta)(x\xi'+y\eta'+z\zeta')}{r^4}
\]
also die eine partielle Kraft
\[
\begin{aligned}
&\frac{x}{r^3}(\xi'\mathrm dx+\eta'\mathrm dy+\zeta'\mathrm dz)-\frac{3x}{2r^4}(x\xi'+y\eta'+z\zeta')\mathrm dr\\[5pt]
&=\tfrac{1}{2}\xi'\mathrm d\frac{xx}{r^3}+\tfrac{1}{2}\eta'\mathrm d\frac{xy}{r^3}+\tfrac{1}{2}\zeta'\mathrm d\frac{xz}{r^3}\\
&\qquad+\frac{\eta'}{2r^3}(x\mathrm dy-y\mathrm dx)+\frac{\zeta'}{2r^3}(x\mathrm dz-z\mathrm dx)
\end{aligned}
\]
\SourceLeaf{625}oder da man vollständige Differentiale weglassen kann
\[
\frac{\eta'(x\eta-y\xi)+\zeta'(x\zeta-z\xi)}{2r^3}=\frac{x(\eta\eta'+\zeta\zeta')-y\xi\eta'-z\xi\zeta'}{2r^3}
\]
Hier ist es nun erlaubt noch zuzusetzen oder wegzulassen
\[
\frac{x\xi\xi'+y\eta\xi'+z\zeta\xi'}{2r^3}
\]
wodurch die Formel symmetrisch in Beziehung auf beide Elemente wird. Wählen wir das letztere, so haben wir die Kraft
\[
\frac{x(-\xi\xi'+\eta\eta'+\zeta\zeta')-y(\xi\eta'+\eta\xi')-z(\xi\zeta'+\zeta\xi')}{2r^3}
\]
Dies erklärt sich durch eine Kraft, die von den \textit{relativen} Bewegungen $\alpha,\beta,\gamma$ abhangt
\[
=-\frac{x(\alpha\alpha+\beta\beta+\gamma\gamma)-2\alpha(x\alpha+y\beta+z\gamma)}{2r^3}
\]

Hier reichen wir nun mit Einer Zusatzkraft aus, die nach $\mathrm{eb}$ mit der Stärke $\frac{vv}{kk}ee$ wirkt, wenn $\mathrm{ea}$ die Richtung der relativen Bewegung $v$ ist
\[
\mathrm{aeb}=180^0-\mathrm{aeE},\qquad\mathrm{eE}\text{ positiv}
\]
\begin{center}\resizebox{.64\textwidth}{!}{\begin{tikzpicture}[x=1bp,y=-1bp,line cap=round,line join=round]
\path[use as bounding box] (0,0) rectangle (890,375);
\begin{pgfinterruptboundingbox}
\draw[line width=2bp] (65,284) -- (618,285);
\draw[line width=2.3bp] (499,360) -- (845,146);
\draw[line width=2.3bp] (822,150) -- (845,146) -- (833,165);
\draw[line width=2.3bp,dash pattern=on 13bp off 8bp] (622,276) -- (749,62);
\draw[line width=2.3bp] (731,78) -- (749,62) -- (746,86);
\draw[fill=white,line width=2bp] (65,284) circle[radius=6bp];
\draw[fill=white,line width=2bp] (618,285) circle[radius=7bp];
\node[inner sep=0pt,outer sep=0pt] at (38,270) {\scalebox{4}{\ensuremath{E}}};
\node[inner sep=0pt,outer sep=0pt] at (622,311) {\scalebox{4}{\ensuremath{e}}};
\node[inner sep=0pt,outer sep=0pt] at (844,130) {\scalebox{4}{\ensuremath{a}}};
\node[inner sep=0pt,outer sep=0pt] at (753,44) {\scalebox{4}{\ensuremath{b}}};
\end{pgfinterruptboundingbox}
\end{tikzpicture}
}\end{center}

Noch einfacher und ganz allgemein wird das Gesetz folgendermaassen ausgedrückt
\begin{center}\resizebox{.64\textwidth}{!}{\begin{tikzpicture}[x=1bp,y=-1bp,line cap=round,line join=round]
\path[use as bounding box] (0,0) rectangle (940,370);
\begin{pgfinterruptboundingbox}
\draw[line width=2.5bp] (70,180) -- (500,42);
\draw[line width=2.5bp] (475,42) -- (500,42) -- (480,57);
\draw[line width=2.7bp] (70,180) -- (866,165);
\draw[line width=2.2bp,dash pattern=on 14bp off 9bp] (82,188) -- (504,304);
\draw[line width=2.3bp] (483,289) -- (504,304) -- (478,305);
\draw[fill=white,line width=2.3bp] (70,180) circle[radius=7bp];
\draw[fill=white,line width=1.8bp] (85,184) circle[radius=4.7bp];
\node[inner sep=0pt,outer sep=0pt] at (45,166) {\scalebox{4}{\ensuremath{E}}};
\node[inner sep=0pt,outer sep=0pt] at (493,21) {\scalebox{4}{\ensuremath{A}}};
\node[inner sep=0pt,outer sep=0pt] at (904,172) {\scalebox{4}{\ensuremath{B}}};
\node[inner sep=0pt,outer sep=0pt] at (493,330) {\scalebox{4}{\ensuremath{C}}};
\end{pgfinterruptboundingbox}
\end{tikzpicture}
}\end{center}

Wenn ein electrisches Element $\mathrm E$ durch die Wirkung eines andern nach der Richtung $\mathrm{EA}$ mit der Kraft $p$ sollicitirt wird, insofern beide in gegenseitiger Ruhe sind, so kommt, im Fall einer gegenseitigen Bewegung, deren Rich%
\SourceLeaf{626}%
tung in Beziehung auf $\mathrm E$, die Gerade $\mathrm{EB}$ und Geschwindigkeit $=v$ ist, zu jener Kraft noch eine zweite hinzu, deren
\[
\text{Stärke}\quad=\frac{pvv}{kk},\qquad\text{Richtung}\quad=\mathrm{EC}
\]
wobei $\mathrm{EA,EB,EC}$ in Einem Planum liegen und $\mathrm{EB}$ mitten zwischen $\mathrm{EA}$ und $\mathrm{EC}$.

\begin{center}[15.]\par\medskip\textit{Kugelfläche.}\end{center}
Es seien $x,y,z$ die Coordinaten eines Punktes in einer auf der Kugelfläche liegenden in sich selbst zurückkehrenden Linie. $xx+yy+zz=rr$

Das Integral
\[
\int\frac{z(x\mathrm dy-y\mathrm dx)}{r(xx+yy)}=V
\]
wird durch die Länge der ganzen Linie genommen verstanden.

Jene Linie scheidet die Kugelfläche in zwei oder mehrere Theile $A,B,C$ u.\ s.\ w.

Es wird dann
\[
V=\alpha A+\beta B+\gamma C+\,.\,.
\]
sein, wo die Coefficienten $\alpha,\beta,\gamma,\,.\,.$ folgenden zwei Bedingungen Genüge leisten:

1.\quad Die Coefficienten zweier an einander grenzenden Stücke sind immer um eine Einheit verschieden, und zwar gehört demjenigen Stück der kleinere Coefficient an, welches gegen die Scheidungslinie ebenso liegt, wie der positive Pol der $z$ gegen den grössten Kreis, der vom Pole der $x$ nach dem Pole der $y$ gezogen ist.

2.\quad Die Summe der Coefficienten, welche denjenigen Stücken angehören, in denen die beiden Pole der $z$ liegen, ist $=0$.

Der Beweis ist leicht geführt, indem man vom negativen zum positiven Pole der $z$ eine unendlich grosse Menge von Halbkreisen zieht.

Der Fall, wo einer der beiden Pole in die Linie selbst fiele, ist durch die Natur des Integrals von selbst ausgeschlossen.

\SourceLeaf{627}
\begin{center}[16.]\par\medskip\textit{Electromotorische Kraft, durch Entstehung eines Stromes.}\end{center}
\begin{center}\resizebox{.76\textwidth}{!}{\begin{tikzpicture}[x=1bp,y=-1bp,line cap=round,line join=round]
\path[use as bounding box] (0,0) rectangle (1130,250);
\begin{pgfinterruptboundingbox}
\draw[line width=2bp,dash pattern=on 11bp off 8bp] (46,79) -- (46,184);
\draw[line width=2bp,dash pattern=on 12bp off 8bp] (46,184) -- (285,185);
\draw[line width=2.8bp] (285,185) -- (1071,186);
\fill (46,70) circle[radius=4bp];
\node[inner sep=0pt,outer sep=0pt] at (47,39) {\scalebox{4}{\ensuremath{P}}};
\node[inner sep=0pt,outer sep=0pt] at (50,211) {\scalebox{4}{\ensuremath{O}}};
\node[inner sep=0pt,outer sep=0pt] at (285,211) {\scalebox{4}{\ensuremath{A}}};
\node[inner sep=0pt,outer sep=0pt] at (1096,174) {\scalebox{4}{\ensuremath{B}}};
\end{pgfinterruptboundingbox}
\end{tikzpicture}
}\end{center}
Electromotorische Kraft in $P$, vermöge Entstehung des Stromes in $A\, .\, .\, .\, B$
\[
=\log\frac{x'+\BareRoot{(x'x'+yy)}}{x^0+\BareRoot{(x^0x^0+yy)}},\quad\text{wenn}\quad0A=x^0,\quad0B=x',\quad0P=y
\]

Electromotorische Kraft in $P$ vermöge Entstehung eines Stromes durch den Kreis
\begin{center}\resizebox{.36\textwidth}{!}{\begin{tikzpicture}[x=1bp,y=-1bp,line cap=round,line join=round]
\path[use as bounding box] (0,0) rectangle (585,510);
\begin{pgfinterruptboundingbox}
\draw[line width=2.5bp] (314,36) .. controls (431.15,36) and (527,131.85) .. (527,249) .. controls (527,366.15) and (431.15,462) .. (314,462) .. controls (196.85,462) and (101,366.15) .. (101,249) .. controls (101,131.85) and (196.85,36) .. (314,36);
\draw[line width=1.9bp,dash pattern=on 10bp off 7bp] (64,250) -- (527,246);
\draw[line width=1.9bp,dash pattern=on 10bp off 7bp] (147,116) -- (319,248);
\fill (64,250) circle[radius=4bp];
\node[inner sep=0pt,outer sep=0pt] at (44,235) {\scalebox{3.7}{\ensuremath{P}}};
\node[inner sep=0pt,outer sep=0pt] at (87,270) {\scalebox{3.7}{\ensuremath{a}}};
\node[inner sep=0pt,outer sep=0pt] at (125,285) {\scalebox{3.7}{\ensuremath{A}}};
\node[inner sep=0pt,outer sep=0pt] at (254,273) {\scalebox{3.7}{\ensuremath{r}}};
\node[inner sep=0pt,outer sep=0pt] at (244,216) {\scalebox{3.7}{\ensuremath{\varphi}}};
\node[inner sep=0pt,outer sep=0pt] at (127,101) {\scalebox{3.7}{\ensuremath{B}}};
\node[inner sep=0pt,outer sep=0pt] at (547,245) {\scalebox{3.7}{\ensuremath{c}}};
\end{pgfinterruptboundingbox}
\end{tikzpicture}
}\end{center}
\[
\begin{aligned}
&\text{1) durch das Stück }AB\,.\,.\text{ proxime }\log\frac{2r\sin\varphi}{a}\\
&\text{2) durch das Stück }BC\text{ proxime }-2-\log\operatorname{tg}\tfrac14\varphi\\
&\text{Zusammen}\quad\log\frac{8r\varphi}{a}-2=\log\frac{8r\varphi}{aee}\\
&\text{Durch den ganzen Kreis}\quad2\log\frac{8r}{aee}
\end{aligned}
\]

\begin{center}[17.]\end{center}
Es entstehe durch Umdrehung eines Kreises, dessen Halbmesser $=\rho$ um eine in jener Ebene liegende Axe, von welcher der Mittelpunkt des Kreises die Entfernung $=R$ hat, ein ringförmiger körperlicher Raum der gleichförmig mit einer Rheophorkette ausgefüllt ist. Die Anzahl der Umwindungen sei $M$. Es wird angenommen, dass $\rho$ gegen $R$ sehr klein sei. Man verlangt die electromotorische Wirkung des Ringes auf einen Punkt, der entweder innerhalb oder sehr nahe am Ringe liegt.
\begin{center}\resizebox{.44\textwidth}{!}{\begin{tikzpicture}[x=1bp,y=-1bp,line cap=round,line join=round]
\path[use as bounding box] (0,0) rectangle (735,410);
\begin{pgfinterruptboundingbox}
\draw[line width=2.4bp] (164,78) .. controls (238.80,78) and (300,139.20) .. (300,214) .. controls (300,288.80) and (238.80,350) .. (164,350) .. controls (89.20,350) and (28,288.80) .. (28,214) .. controls (28,139.20) and (89.20,78) .. (164,78);
\draw[line width=1.9bp,dash pattern=on 11bp off 8bp] (165,213) -- (687,209);
\draw[line width=1.9bp,dash pattern=on 11bp off 8bp] (685,24) -- (693,374);
\end{pgfinterruptboundingbox}
\end{tikzpicture}
}\end{center}

Es sei $a$ die Distanz des Punktes von der Centrallinie des Ringes, so wird die verlangte Wirkung sein
\[
\begin{aligned}
2M\log\frac{8R}{aee}&=2M\left\{\log\frac{8R}{a}-2\right\}\quad\text{wenn}\quad a>\rho\\
2M\left(\log\frac{8R}{\rho}-\tfrac32-\frac{aa}{2\rho\rho}\right)&\quad\text{wenn}\quad a<\rho
\end{aligned}
\]

\SourceLeaf{628}
Der mittlere Werth für alle \textit{im} Ringe gleichförmig vertheilten Punkte ist
\[
2M\left\{\log\frac{8R}{\rho}-\tfrac74\right\}
\]
oder die ganze electromotorische Kraft, welche eindem inducirenden Ringe gleichförmig eingewirkter von $m$ Umwindungen erleidet
\[
=4\pi mMR\left\{\log\frac{8R}{\rho}-\tfrac74\right\}=E
\]

Bei gegebenem Drahtvorrath für jede Kette ist $mR,MR$ und $\dfrac{m}{\rho\rho}$ gegeben, da nun
\[
E=16\pi.\ MR.\ (mR)^{\frac23}.\ \left(\frac{m}{\rho\rho}\right)^{\frac13}.\ \left(\frac{8R}{\rho}\right)^{-\frac23}\left(\log\frac{8R}{\rho}-\tfrac74\right)
\]
so muss, damit $E$ \textit{ein Maximum} werde
\[
\left(\frac{8R}{\rho}\right)^{-\frac23}\left(\log\frac{8R}{\rho}-\tfrac74\right)
\]
ein Maximum werden, oder $\left(\dfrac{8R}{\rho}\right)^{\frac23}=x$ gesetzt, muss
\[
\frac{\log x^{\frac32}-\tfrac74}{x}=\tfrac32\left\{\frac{\log x-\tfrac76}{x}\right\}
\]
ein Maximum werden.

Dies geschieht, wenn $\log x=\tfrac{13}{6}$
\[
\begin{aligned}
&\text{oder}\quad\log\frac{8R}{\rho}=\tfrac{13}{4}\quad\text{wird}\\
&\text{oder}\quad\log\operatorname{Brigg}\frac{8R}{\rho}=1{,}4114580\\
&\text{oder}\quad\frac{8R}{\rho}=25{,}79\\
&\hphantom{\text{oder}\quad}R=3{,}22\rho
\end{aligned}
\]

\begin{center}[18.]\end{center}
Wenn man zu \textit{den Massen im Innern eines körperlichen Raumes} noch die ihnen für den äussern Raum aequivalirenden auf der Oberfläche mit entgegengesetzten Zeichen beifügt, so erhält man einen Körper als Träger von positiven und negativen Massen, deren Complex auf alle Punkte des äussern Raumes gar keine Anziehungskraft ausübt.

\SourceLeaf{629}
Man beweiset leicht

1.\quad dass in Folge der Reaction äusserer Massen jener Körper auch im Gleichgewicht bleibt

2.\quad dass der Körper, wenn die betreffenden Massen magnetische Fluida sind, auch auf einen Rheophor gar keine Kraft ausübt.

Schwerer aber

3.\quad dass auch trotz der Reaction des Rheophors jener Körper im Gleichgewichte bleibt.

Das letzte beruht auf folgenden Momenten:

Es sei $\mathrm dm$ ein Element des Körpers; $x,y,z$ dessen Coordinaten; die Characteristik $\int$ beziehe sich auf alle $\mathrm dm$. Es sei ferner $\mathrm ds$ ein Element eines galvanischen Stroms, $a,b,c;\ a+\mathrm da,\ b+\mathrm db,\ c+\mathrm dc$ die Coordinaten seiner Endpunkte, wo also $a,b,c$ Functionen von $s$. Die Intensität des Stromes $=1$. Charakteristik $S$ Summation in Beziehung auf $\mathrm ds$.

Sind $X,Y,Z$ die Componenten der ganzen auf $\mathrm dm$ wirkenden beschleunigenden Kraft, so sind die Bedingungen des Gleichgewichts bekanntlich
\[
\begin{gathered}
\int X\mathrm dm,\int Y\mathrm dm,\int Z\mathrm dm,\int(zY-yZ)\mathrm dm,\\
\int(xZ-zX)\mathrm dm,\int(yX-xY)\mathrm dm\text{ alle }=0
\end{gathered}
\]

Es ist $(x-a)^2+(y-b)^2+(z-c)^2=rr$ gesetzt (die positiven $x,y,z$ bez. nach vorn, rechts oben gerichtet)
\[
\begin{aligned}
X&=S\frac{(y-b)\mathrm dc-(z-c)\mathrm db}{r^3},\\
Y&=S\frac{(z-c)\mathrm da-(x-a)\mathrm dc}{r^3},\\
Z&=S\frac{(x-a)\mathrm db-(y-b)\mathrm da}{r^3}
\end{aligned}
\]
\[
\text{also}\qquad\int X\mathrm dm=S(\eta\mathrm dc-\zeta\mathrm db)
\]
\[
\begin{aligned}
&\int(yX-xY)\mathrm dm\\
&=SV\mathrm dc+S\{(\xi c-\zeta a)\mathrm da+(\eta c-\zeta b)\mathrm db+(a\xi+b\eta+c\zeta)\mathrm dc\}\\
&\qquad-\int z\mathrm dmS\frac{\mathrm d\frac1r}{\mathrm ds}\mathrm ds
\end{aligned}
\]
[wo $\displaystyle\int\frac{\mathrm dm}{r}=V,\quad\int\frac{x-a}{r^3}\mathrm dm=\xi$ u.\ s.\ f. gesetzt sind, welche alle zu Null werden.]

\medskip
1836 Februar 18.

Viel einfacher wird die Ableitung auf folgende Art gemacht.

Wenn $x,y,z;\ x+\mathrm dx,\ y+\mathrm dy,\ z+\mathrm dz$ die Coordinaten zweier beliebiger einander unendlich naher Punkte sind, so muss die Variation von
\[
\mathrm dx^2+\mathrm dy^2+\mathrm dz^2
\]
\SourceLeaf{630}
unabhängig von den Werthen von $\mathrm dx,\mathrm dy,\mathrm dz,=0$ werden. Zur Abkürzung bezeichnen wir $\delta x,\delta y,\delta z$ mit $\alpha\xi,\alpha\eta,\alpha\zeta$, wo $\alpha$ einen constanten unendlich kleinen Coëfficienten bedeutet. Man hat also
\[
\begin{aligned}
\delta(\mathrm dx^2+\mathrm dy^2+\mathrm dz^2)&=2\alpha(\mathrm dx.\mathrm d\xi+\mathrm dy.\mathrm d\eta+\mathrm dz.\mathrm d\zeta)\\
&=2\alpha\Big\{\frac{\mathrm d\xi}{\mathrm dx}.\mathrm dx^2+\frac{\mathrm d\eta}{\mathrm dy}.\mathrm dy^2+\frac{\mathrm d\zeta}{\mathrm dz}.\mathrm dz^2\\
&\quad+\left(\frac{\mathrm d\xi}{\mathrm dy}+\frac{\mathrm d\eta}{\mathrm dx}\right)\mathrm dx\mathrm dy+\left(\frac{\mathrm d\xi}{\mathrm dz}+\frac{\mathrm d\zeta}{\mathrm dx}\right)\mathrm dx\mathrm dz\\
&\quad+\left(\frac{\mathrm d\eta}{\mathrm dz}+\frac{\mathrm dz}{\mathrm dy}\right)\mathrm dy\mathrm dz\Big\}
\end{aligned}
\]
Offenbar muss also sein
\[
\begin{aligned}
1.\quad0&=\frac{\mathrm d\xi}{\mathrm dx}\\
2.\quad0&=\frac{\mathrm d\eta}{\mathrm dy}\\
3.\quad0&=\frac{\mathrm d\zeta}{\mathrm dz}\\
4.\quad\frac{\mathrm d\xi}{\mathrm dy}&=-\frac{\mathrm d\eta}{\mathrm dx}=r\\
5.\quad\frac{\mathrm d\zeta}{\mathrm dx}&=-\frac{\mathrm d\xi}{\mathrm dz}=q\\
6.\quad\frac{\mathrm d\eta}{\mathrm dz}&=-\frac{\mathrm d\zeta}{\mathrm dy}=p
\end{aligned}
\]
Man hat aus (6) (4) (5) (6), aus (6) (2), aus (6) (3):
\[
\begin{aligned}
\frac{\mathrm dp}{\mathrm dx}&=\frac{\mathrm dd\eta}{\mathrm dx\mathrm dz}=-\frac{\mathrm dd\xi}{\mathrm dz\mathrm dy}=\frac{\mathrm dd\zeta}{\mathrm dy\mathrm dx}=-\frac{\mathrm dp}{\mathrm dx}=0\\
\frac{\mathrm dp}{\mathrm dy}&=\frac{\mathrm dd\eta}{\mathrm dy\mathrm dz}=0\\
\frac{\mathrm dp}{\mathrm dz}&=-\frac{\mathrm dd\zeta}{\mathrm dz\mathrm dy}=0
\end{aligned}
\]
also $p$ constant,

ebenso folgt $q,r$ constant

und daher
\[
\begin{aligned}
\xi&=a+ry-qz\\
\eta&=b+pz-rx\\
\zeta&=c+qx-py
\end{aligned}
\]

\begin{center}[19.]\par\medskip\textit{Beweis von} \textsc{Ampères} \textit{Fundamentalsatze.}\end{center}
Über eine begrenzte Fläche (I), in der jedem unbestimmten Punkte die Coordinaten $x,y,z$ angehören, sei positives magnetisches Fluidum gleichförmig so verbreitet, dass auf die Flächeneinheit das Quantum magnetischen Fluidums $=k$ komme.

Ein Element eines galvanischen Stroms von der Intensität $i$ erstrecke sich von $0,0,0$, bis $0,0,\zeta$.

\SourceLeaf{631}
Zur Abkürzung schreibe man
\[
r=\BareRoot{(xx+yy+zz)}
\]

Einem Elemente $\omega$ der Coordinatenebene der $x,y$ entspricht das Element der Fläche (I) $\, .\, .\, .\,\omega\BareRoot{\left(1+\left(\dfrac{\mathrm dz}{\mathrm dx}\right)^2+\left(\dfrac{\mathrm dz}{\mathrm dy}\right)^2\right)}$ und das Quantum magnetischen Fluidums $k\omega\BareRoot{\left(1+\left(\dfrac{\mathrm dz}{\mathrm dx}\right)^2+\left(\dfrac{\mathrm dz}{\mathrm dy}\right)^2\right)}$. Dessen Wirkung auf das Stromelement $i\zeta$ zerlegt sich also in die drei partiellen Kräfte
\[
\begin{gathered}
ik\zeta\omega.\frac{y}{r^3}\BareRoot{\left(1+\left(\frac{\mathrm dz}{\mathrm dx}\right)^2+\left(\frac{\mathrm dz}{\mathrm dy}\right)^2\right)}\\
-ik\zeta\omega.\frac{x}{r^3}\BareRoot{\left(1+\left(\frac{\mathrm dz}{\mathrm dx}\right)^2+\left(\frac{\mathrm dz}{\mathrm dy}\right)^2\right)}\\
0
\end{gathered}
\]

Nehmen wir jetzt eine zweite Fläche (II) mit (I) parallel in der unendlich kleinen Entfernung $\varepsilon$ unter dieser, d.\ i. jedem Punkte $x,y,z$, in (I) entspreche in II der Punkt
\[
\begin{gathered}
x+\frac{\varepsilon}{\BareRoot{\left(1+\left(\frac{\mathrm dz}{\mathrm dx}\right)^2+\left(\frac{\mathrm dz}{\mathrm dy}\right)^2\right)}}.\left(\frac{\mathrm dz}{\mathrm dx}\right)\\
y+\frac{\varepsilon}{\BareRoot{\left(1+\left(\frac{\mathrm dz}{\mathrm dx}\right)^2+\left(\frac{\mathrm dz}{\mathrm dy}\right)^2\right)}}.\left(\frac{\mathrm dz}{\mathrm dy}\right)\\
z-\frac{\varepsilon}{\BareRoot{\left(1+\left(\frac{\mathrm dz}{\mathrm dx}\right)^2+\left(\frac{\mathrm dz}{\mathrm dy}\right)^2\right)}}
\end{gathered}
\]

Über diese zweite Fläche sei negatives magnetisches Fluidum dergestalt verbreitet, dass jeder Flächentheil von II eben so viel negatives Fluidum enthalte, als der entsprechende Flächentheil von I positives. Die Gesammtwirkung derjenigen Fluida, die auf den einander entsprechenden Elementen von I. II enthalten sind, werden demnach
\[
\begin{aligned}
\frac{\varepsilon ik\zeta.\omega}{r^5}\left\{3xy.\left(\frac{\mathrm dz}{\mathrm dx}\right)+(3yy-rr)\left(\frac{\mathrm dz}{\mathrm dy}\right)-3yz\right\}&=X\omega\\
-\frac{\varepsilon ik\zeta.\omega}{r^5}\left\{(3xx-rr)\left(\frac{\mathrm dz}{\mathrm dx}\right)+3xy\left(\frac{\mathrm dz}{\mathrm dy}\right)-3xz\right\}&=Y\omega\\
0&=Z\omega
\end{aligned}
\]

Fangen wir mit der Umformung von $\omega X$ an, welches wir zuerst in die Form setzen
